\section{Алгоритмы начального снимка и потоковой актуализации}
\label{sec:snapshot-stream}

\subsection{Начальный снимок как этап инициализации}

При запуске сервис запрашивает полный снимок таблицы маршрутизации
отдельно для IPv4 и IPv6. Такой разделённый подход упрощает логическую
обработку, последующее распараллеливание и раздельную диагностику двух
семейств адресов.

После получения маршрутов из GoBGP данные ещё не являются готовой
таблицей \textit{prefix~$\rightarrow$~ASN}. Исходное множество маршрутов
может содержать перекрывающиеся префиксы, а потребителям сервиса
требуется представление, в котором для адресной области определён
наблюдаемый ASN.

\subsection{Разрешение перекрывающихся префиксов}

Одной из важных задач при построении полной выгрузки является устранение
неоднозначностей, связанных с перекрывающимися сетевыми префиксами. Если
непосредственно скопировать все маршруты в хранилище, внешний
пользователь при построении полной таблицы получит набор перекрывающихся
диапазонов, неудобный для клиентских сценариев анализа.

Поэтому между этапом чтения маршрутов и этапом загрузки во внутреннее
хранилище выполняется нормализация диапазонов: исходный набор маршрутов
приводится к непересекающемуся множеству диапазонов, в котором каждой
области сопоставляется доминирующий ASN. На уровне идеи это близко к
обработке интервалов сканирующей прямой: события «начало диапазона» и
«конец диапазона» обрабатываются в порядке адресов, а в каждой точке
поддерживается множество активных префиксов с выбором наиболее
специфичного из них.

Этот этап особенно важен для полной выгрузки: клиенты получают не
произвольный внутренний набор префиксов, а нормализованное представление
состояния.

\subsection{Замена полного состояния}

После построения снимка сервис выполняет операцию \texttt{Replace} для
выбранного внутреннего хранилища. Для прежнего варианта на
\texttt{ranger} это естественная операция, поскольку вся модель хранения
основана на полной замене. Для ART-режима \texttt{Replace} тоже
допустим, но используется только на этапе начальной инициализации и
ресинхронизации, а не как нормальный способ обработки каждого
update-события.

\subsection{Потоковая обработка изменений}

После начальной загрузки сервис подписывается на поток BGP-событий.
Каждое событие интерпретируется как одна из двух базовых операций над
\texttt{prefixStore}:

\begin{itemize}
\item \texttt{Upsert} маршрута;
\item \texttt{Delete} маршрута.
\end{itemize}

В потоковой модели именно эти операции становятся главным рабочим
режимом. Их обработка должна:

\begin{enumerate}
\item быть локальной, а не глобальной;
\item изменять только затронутый участок состояния;
\item сохранять корректность \texttt{Lookup};
\item порождать уведомления для подписчиков через подсистему рассылки.
\end{enumerate}

\subsection{Пересинхронизация после ошибок}

Потоковая модель имеет естественное ограничение: если соединение
оборвалось, нельзя безусловно предположить, что локальное состояние
осталось согласованным с источником. Поэтому система использует стратегию
повторного подключения и ресинхронизации, изображённую как конечный
автомат на рис.~\ref{fig:lifecycle}:

\begin{enumerate}
\item поток завершается или переходит в ошибку;
\item клиент закрывает соединение;
\item заново считывается полный снимок;
\item выполняется \texttt{Replace} для всего хранилища;
\item подписка на поток событий открывается повторно.
\end{enumerate}

С точки зрения модели состояния это соответствует идее Route
Refresh~\cite{rfc2918,rfc7313} и подходу «снимок как начальная точка» в
системах обработки потоков~\cite{akidau2015dataflow}: поток трактуется
как эволюция состояния относительно базовой точки, и если эта точка
стала недостоверной, она строится заново.

\begin{figure}[ht]
\centering
\begin{tikzpicture}[node distance=0.7cm and 1.0cm,every node/.style={font=\footnotesize}]
\node[state, fill=gray!10] (init)    {Init};
\node[state, fill=gray!10, right=of init]   (conn)    {Connect};
\node[state, fill=gray!10, right=of conn]   (snap)    {Snapshot\\\texttt{ListPath}};
\node[state, fill=gray!10, below=of snap]   (load)    {Normalize\\\& Load};
\node[state, fill=gray!10, left=of load]    (watch)   {Watch loop\\\texttt{WatchEvent}};
\node[state, fill=red!10, below=of watch]   (err)     {Error /\\disconnect};

\draw[msarr] (init) -- (conn);
\draw[msarr] (conn) -- (snap);
\draw[msarr] (snap) -- (load);
\draw[msarr] (load) -- (watch);
\draw[msarr] (watch) -- (err);
\draw[msarr] (err.west) -- ++(-1,0) |- (conn.south);
\end{tikzpicture}
\caption{Конечный автомат жизненного цикла повторного подключения и ресинхронизации}
\label{fig:lifecycle}
\end{figure}

\subsection{Модель состояния и корректность}

Снимок и поток удобно рассматривать не как противоположности, а как две
фазы жизни одного и того же материализованного состояния. Снимок нужен
при старте, при пересинхронизации и при выгрузке полной таблицы наружу.
Поток снижает задержку актуализации состояния, применяет
инкрементальные изменения отдельных префиксных записей и передаёт
обновления подписчикам.

Порядок применения событий задаётся самой сессией с GoBGP: сервис
обрабатывает их в том порядке, в котором они приходят. По терминологии
работ по обработке потоков~\cite{akidau2015dataflow} это
\textit{processing-time}, а не \textit{event-time}. Для прикладной
задачи такой выбор обоснован по двум причинам. Во-первых, истинное время
события у исходного BGP-роутера сервису недоступно -- между ним и
сервисом стоит GoBGP, агрегирующий поток. Во-вторых, важна не глобальная
синхронизация по астрономическим часам, а локальная согласованность
шагов: после применения события $e_i$ и до применения $e_{i+1}$ ответы
на \texttt{Lookup} должны соответствовать состоянию после $e_i$. Эта
гарантия обеспечивается дисциплиной работы с \texttt{prefixStore}:
события обрабатываются по одному, а видимость их результата управляется
блокировками внутри шардов.

Кратко состояние таблицы записывается как
\[
S_n = f(S_0, e_1, e_2, \ldots, e_n),
\]
где $S_0$ -- результат начальной загрузки, а $e_1, \ldots, e_n$ --
последовательно применённые BGP-события. Любой ответ \texttt{Lookup}
клиента $C$ в момент $t_C$ соответствует состоянию после некоторого числа
применённых событий $S_k$, $k \le n(t_C)$. Полная выгрузка
\texttt{GetPrefixesASTable} возвращает текущее представление состояния
сервиса; в шардированной реализации её согласованность уточняется в
главе~\ref{sec:art-storage}. Подписка \texttt{SubscribePrefixASUpdates}
даёт поток дельт, который позволяет клиенту, начав с $S_k$, восстановить
$S_{k+1}, S_{k+2}, \ldots$ Из этой записи видно главное: снимок и поток
связаны единой семантикой, а \texttt{Replace} после ресинхронизации --
переустановка базовой точки $S_0$ при сохранении того же типа функции
$f$.

%==============================================================
